% !TEX root = ../notes.tex

\documentclass[../notes.tex]{subfiles}

\begin{document}

\section{April 29}
Here we go.

\subsection{Cluster Varieties}
For today, we will choose a finite type acyclic cluster algebra $\mc A$, for which we may assume has initial seed $(x,B)$ with $B$ whose lower-triangular submatrix is given by the associated Cartan matrix.
\begin{definition}[cluster variety]
	An \textit{acyclic cluster variety} is an affine variety isomorphic to $\Spec\mc A$ for some acyclic cluster algebra $\mc A$. In the sequel, we will abbreviate this to a \textit{cluster variety}.
\end{definition}
\begin{remark} \label{rem:general-cluster-variety}
	There is a more general definition due to Fock--Goncharov and Gross--Hacking--Keel: namely, one can define a cluster variety to be covered by some algebraic tori $T_\Sigma$, which we glue together via some exchange relations. For example, there exist projective cluster varieties.
\end{remark}
It will be helpful to have a version of a cluster variety which is manifestly integral and finite type.
\begin{definition}[bounded cluster algebra]
	Fix a cluster algebra $\mc A$ with initial seed $(x,B)$. Then the \textit{lower bound cluster algebra} $\mc L$ of $\mc A$ is the subalgebra generated over $\ZZ\left[x_1,\ldots,x_n\right]$ by the cluster variables which are one mutation away. Similarly, we define the \textit{upper bound cluster algebra} $\mc U$ of $\mc A$ is the intersection
	\[\mc U\coloneqq\bigcap_{z}\ZZ\left[z_1^\pm,\ldots,z_n^\pm\right],\]
	where $z$ varies over all seeds.
\end{definition}
\begin{remark}
	Note that $\mc L\subseteq\mc A\subseteq\mc U$. Indeed, $\mc L\subseteq\mc A$ has no content, and $\mc A\subseteq\mc U$ follows from \Cref{thm:laurent}.
\end{remark}
\begin{remark}
	One can glue the tori $\Spec\ZZ\left[z_1^\pm,\ldots,z_n^\pm\right]$ together to produce $\mc A$ in some cases. This explains how \Cref{rem:general-cluster-variety} generalizes our definition of a cluster variety.
\end{remark}
In fact, in the acyclic case, these bounds collapse.
\begin{theorem}[Berenstein--Fomin--Zelevinsky]
	If $\mc A$ is an acyclic cluster algebra, then $\mc L=\mc A=\mc U$.
\end{theorem}
\begin{proof}[Sketch]
	One shows $\mc U\subseteq\mc L$ using the intersection lemma discussed in \Cref{thm:laurent}.
\end{proof}
\begin{example}
	From \Cref{ex:clusterl-alg-a3}, one has that $\mc A=\ZZ[x_1,x_2,(x_2+1)/x_1,(x_1+1)/x_2]$. Indeed, the last cluster variable is $(x_1+x_2+1)/(x_1x_2)$, which is $(x_2+1)/x_1\cdot(x_1+1)/x_2-1$.
\end{example}
For example, one can now say something about friezes.
\begin{definition}
	Fix an $n\times n$ generalized Cartan matrix. Then for each $i$, define the polynomial
	\[f_{C,i}\coloneqq x_iy_i-\prod_{j<o}x_j^{\left|c_{ji}\right|}-\prod_{j>i}x_j^{\left|c_{ji}\right|}\]
	in $\ZZ[x_1,\ldots,x_n,y_1,\ldots,y_n]$. We now define $Y_1^{\mathrm{frieze}}(C)$ to be the $n$-dimensional subvariety cut out by these functions. We may abbreviate $Y_1^{\mathrm{frieze}}(C)$ to $Y_1^{\mathrm{frieze}}(\Delta)$, where $\Delta$ denotes the generalized Dynkin type.
\end{definition}
\begin{remark}
	The notation $Y_1$ is intended to be similar to the notation of a modular curve.
\end{remark}
\begin{remark}
	If $\mc A$ is an acyclic cluster algebra, then it turns out that $\Spec\mc A\cong Y_1^{\mathrm{frieze}}(C)$.
\end{remark}
In fact, the remark upgrades as follows.
\begin{theorem}[de Saint Germain--Huang--Lin]
	Fix a generalized $n\times n$ generalized Cartan matrix $C$. Then there is a bijection between $Y_1^{\mathrm{frieze}}(C)(\NN)$ and $\op{Frieze}(C,\NN)$, which sends the variables $(x_1,\ldots,x_n)$ to the first column of the frieze. Here, $Y_1^{\mathrm{frieze}}(C)(\NN)$ means that all the variables $(x_1,\ldots,x_n,y_1,\ldots,y_n)$ must get sent to positive integers.
\end{theorem}
\begin{remark}
	The values of the other variables $y_1,\ldots,y_n$ can also be seen in other entries of the corresponding frieze
\end{remark}
Thus, $Y_1^{\mathrm{frieze}}(C)$ can be understood as a moduli space of friezes.

\subsection{Diophantine Geometry}
We will get some motivation from the following comparison of trichotomies.
\begin{itemize}
	\item Fano varieties have $-K_X$ ample.
	\item Calabi--Yau varieties have $K_X$ trivial.
	\item General type varieties have $K_X$ big.
\end{itemize}
On the other hand, one has some results for curves.
\begin{itemize}
	\item Fano cures have genus zero, which have a dense subset of points after a field extension.
	\item Calabi--Yau curves have genus one, which are elliptic curves (once attaining a point), which remain hard to understand.
	\item General type curves have genus at least two, which have finitely many points.
\end{itemize}
The first point produces the following conjecture.
\begin{conj}[Campana]
	A Fano variety $X$ over a number field $K$ admits a finite extension $L$ for which $X(L)$ is Zariski dense.
\end{conj}
The last point produces the following conjecture.
\begin{conj}[Bombieri--Lang]
	Fix a variety $X$ of general type over a number field $K$. Then $X(K)$ is not Zariski dense.
\end{conj}
\begin{remark}
	There are similar statements for integral points.
\end{remark}
The Calabi--Yau case remains hard to understand; for example, one can hunt for infinite orbits, but there are other ways to find rational points. For example, there is a conjecture that a K3 surface admitting local points and no Brauer--Manin obstruction admit rational points.

This trichotomy is relevant for us because our cluster varieties turn out to live in the interesting Calabi--Yau case.
\begin{theorem}
	Every cluster variety is log Calabi--Yau.
\end{theorem}
\begin{remark}
	Here, we mean that one can glue the tori together with the ``boundary'' of varieties cut out by $x_i=0$ along $\Omega_\Sigma=dx_1/x_1\land\cdots\land dx_n/x_n$. The adjective ``log'' is present because $\Omega_\Sigma$ is basically the wedge of the $d\log x_i$s.
\end{remark}

\subsection{Counting Friezes}
Note that it is not so hard to find rational points on $\Spec\mc A$, so we will focus on the harder question of finding integral points or even positive integral points, where a positive integral point is one in which the cluster variables go to positive integers.

It will help to have the following automorphism. The following is due to Kontsevich--Soiberlman, Keller, Fock--Goncharov, Gross--Hacking--Keel, and Fomin--Zelevinsky.
\begin{theorem} \label{thm:tau-exists}
	Fix a generalized Dynkin type $\Delta$.
	\begin{listalph}
		\item There is a volume-preserving automorphism $\tau\colon Y_1^{\mathrm{frieze}}(\Delta)\to Y_1^{\mathrm{frieze}}(\Delta)$.
		\item Furthermore, $\tau$ is periodic if and only if $\Delta$ is finite; and if $\Delta$ is finite, then the period divides $h+2$, where $h$ is the Coxeter number of $\Delta$.
	\end{listalph}
\end{theorem}
\begin{remark}
	This $\tau$ is given by horizontally translating a frieze one column. Thus, the existence of the automorphism and the fact that it preserves volumes is not hard to check. The difficulty is the period; for example, for $\mf{sl}_3$, this follows from the work of Gauss.
\end{remark}
\begin{remark}
	This $\tau$ has many names: it can be the Donaldson--Thomas transformation, the maximal green sequence, the Auslander--Reiten translation, or the Fomin--Zelevinsky twist.
\end{remark}
We can use this $\tau$ to understand positive integral points.
\begin{theorem}[Zhang]
	Fix a generalized Dynkin type $\Delta$. Then $Y_1^{\mathrm{frieze}}(\Delta)(\NN)$ is finite if and only if $\Delta$ is finite type, and the values are given by the table in \Cref{thm:count-frieze}.
\end{theorem}
\begin{proof}[Sketch]
	In the case where $\Delta$ is infinite, the result is due to Morier--Genoud. Here are two arguments.
	\begin{itemize}
		\item One can construct many maps $\mc A\to\ZZ$, basically by sending one collection of cluster variables to $1$. There are infinitely many sets of cluster variables, so this produces infinitely many friezes.
		\item Alternatively, as soon as we have a single positive integral point (by sending cluster variables to $1$), one can use \Cref{thm:tau-exists} to show that the orbit of this point has infinite size. Roughly speaking, one knows that $\tau$ preserves positivity and integrality, and one can show that the size of the point grows.
	\end{itemize}
	We now turn to the finite type case. Define $Y_0^{\mathrm{frieze}}(\Delta)$ to be the quotient of $Y_1^{\mathrm{frieze}}(\Delta)$ by $\langle\tau\rangle$, which is a quotient by a finite group because $\tau$ now has finite order. We let $\OO_\tau(x)$ denote the image of $x\in Y_1^{\mathrm{frieze}}(\Delta)$.

	Now, view $\log\OO_\tau(x)$ as a vector and then apply the matrix $C$ to it. It was shown by Muller that $C\cdot\log\OO_\tau(x)$ is in a bounded set. However, Lusztig--Tits showed that $C^{-1}$ has positive rational coefficients if and only if $C$ is finite type, so one finds that the coordinates of $\log\OO_\tau(x)$ are bounded by the coordinates of $C^{-1}(1,\ldots,1)$. This shows that there can only be finitely many orbits $\OO_\tau(x)$, so there are only finitely many points because $\tau$ has finite order!

	It remains to actually count the number of friezes. For example, for $\Delta=E_8$, the bound provided by the previous paragraph is on the order of $10^{50}$ for the height of the orbit, which is much too large. We will do this with the following trick.
	\begin{lemma}
		Fix a finite Dynkin type $\Delta$. Let $i$ be a vertex in $\Delta$ of degree $1$. Then the intersection of $Y_1^{\mathrm{frieze}}(\Delta)$ with $x_i=1$ is isomorphic to $Y_1^{\mathrm{frieze}}(\Delta\setminus\{1\})$.
	\end{lemma}
	\begin{proof}
		One can see this explicitly on the level of the equations. Alternatively, this amounts to saying that a frieze whose bottom row is all $1$s is simply a frieze for the smaller Dynkin type.
	\end{proof}
	Thus, we may assume that all nodes $i$ of degree $1$ have $x_i\ge2$. From here, a long na\"ive calculation over a few days can complete the proof. For example, in $E_8$, there are only four new friezes not from a smaller Dynkin type.
\end{proof}
\begin{remark}
	It is an interesting question to find the number of $\tau$-orbits among the friezes. Professor Zhang has threatened to put this on the homework.
\end{remark}
\begin{example}
	It turns out that $Y_1^{\mathrm{frieze}}\left(\begin{bsmallmatrix}
		2 & -a \\ -b & 2
	\end{bsmallmatrix}\right)$ can be presented with the equation
	\[xyz=\left(x^a+1\right)^b+y.\]
	It follows that there are $5$ points if $ab=1$, $6$ points if $ab=2$, and $9$ points of $ab=9$, and there are infinitely many points if $ab\ge4$. More generally, if $G(x,y)$ has positive integral coefficients, we are able to determine when $xyz=G(x,y)$ should have infinitely many positive integral solutions. For example, there are infinitely many when $\deg_xG$ and $\deg_yG$ are both at least three and has non-vanishing constant term.
\end{example}

\end{document}